% ID: SEL-D0081-OB-E03
% ESTADO: cotejada. Segunda pasada visual y borrador compilado revisados.
% Fuente DOC: Selectividad/sel_2012_matematicas_aplicadas/3 MATEMÁT. APLIC. CC. SS/EXAMEN 2-ANDALUCÍA.doc
% SHA-256 original: bcec57ae525fa5e23ce74631274a2abda0b5b16ac47ce6108ba8b5572d4aa634
% Página convertida: 2; opción B; ejercicio 3.

Se sabe que el $44\%$ de la población activa de cierta provincia está formada por mujeres. También se sabe que, de ellas, el $25\%$ están en paro y que el $20\%$ de los hombres de la población activa también están en paro.
\begin{enumerate}[label=\alph*)]
\item \textbf{(1.25 puntos)} Elegida, al azar, una persona de la población activa de esa provincia, calcule la probabilidad de que esté en paro.
\item \textbf{(1.25 puntos)} Si hemos elegido, al azar, una persona que trabaja, \textquestiondown{}cuál es la probabilidad de que sea hombre?
\end{enumerate}
